\uuid{nSgB}
\titre{Zéros de fonction}
\auteur{bodin ; hanani ; mzari}
\organisation{exo7}
\licence{CC BY-NC-SA 4.0}
\datecreate{2020-01-20}
\niveau{L1}
\module{Analyse}
\chapitre{Continuité, limite et étude de fonctions réelles}
\sousChapitre{Continuité : théorie}
\difficulte{1}
\theme{Théorèmes des valeurs intermédiaires, zéros de fonction}
\exo7qcmid{lille-1.continuite.019}

\contenu{
\begin{question}
\qclassification{123.01,123.02}
\qtags{motcle=zéros de fonction}
Quelles assertions peut-on déduire du théorème des valeurs intermédiaires ?
\begin{answers}
    \good{$\sin(x) - x^2 + 1$ s'annule sur $[0,\pi]$.}

    \good{$x^5-37$ s'annule sur $[2,3]$.}

    \good{$\ln(x+1)-x+1$ s'annule sur $[0,+\infty[$.}

    \bad{$e^x+e^{-x}$ s'annule sur $[-1,1]$.}
\end{answers}
\begin{explanations}
Pour montrer qu'une fonction continue $f$ s'annule sur un intervalle $[a,b]$, il est suffisant de montrer que $f(a)>0$ et $f(b)<0$ (ou l'opposé).
\end{explanations}
\end{question}
}
